% Empirical application from the frozen 2002 pharmaceutical-tablet NIR archive.
\chapter{Empirical Application: Pharmaceutical-Tablet Spectroscopy}
\label{chap:spectroscopy-application}

The controlled experiments in Chapters~\ref{chap:estimation-results} and
\ref{chap:inference-results} separate statistical regularisation from
numerical approximation because the data-generating slope and covariance
operator are known.  This chapter asks whether the resulting mechanisms remain
visible in an observed functional regression problem.  Near-infrared (NIR)
spectra are used to predict the active-ingredient assay of pharmaceutical
tablets with the four complete procedures studied in the thesis: CG--FPLS,
raw FPLS, Arnoldi FPLS, and FPCR.  The application concerns prediction and
numerical stability.  The physical slope is unknown, so it cannot assess slope
recovery or support the oracle-assisted inference studied in Chapters
\ref{chap:inference-methodology}--\ref{chap:inference-results}.

The empirical protocol was frozen before any of the four procedures was fitted
to the supplied test observations.  That test set is nevertheless described
as a \emph{designated benchmark}, not a prospectively blind holdout: its
responses and preliminary ordinary PCR/PLS performance had already been
inspected while selecting a dataset.  This qualification is material to the
scope of the evidence.  The chapter therefore reports a conditional method
comparison on the supplied tablets rather than an estimate of future
deployment performance.

\section{Data and prediction problem}
\label{sec:spectroscopy-data}

The data come from the 2002 pharmaceutical-tablet NIR ``Shoot-out'' described
by \textcite{hopkins_2003} and distributed in MATLAB format by
\textcite{eigenvector_tablets}.  The archive contains spectra from two
spectrometers and three supplied splits: 155 calibration tablets, 40
validation tablets, and 460 test tablets.  These add to 655 tablets, although
the online catalogue reports 654.  All 655 complete archive rows are retained,
and the discrepancy is recorded rather than resolved through an undocumented
deletion.  The two instruments measure the same tablets; they are paired
measurements and are never stacked as independent observations.

The primary predictor is the instrument-1 spectrum at all $T=650$ wavelengths
from 600 to 1898~nm in 2-nm increments.  The sole response is assay, the
archive's third response column.  Its physical unit is not stated in the
distributed object, so prediction errors are reported in the supplied assay
units.  The calibration and validation rows form a 195-tablet development
sample, while the 460 test rows form the designated benchmark.  Figure
\ref{fig:spectroscopy-overview} shows substantial common spectral structure,
a visibly less regular high-wavelength tail, and the unusually narrow assay
range of the supplied validation split.

\begin{figure}[htbp]
    \centering
    \includegraphics[width=0.82\textwidth]{tablet_spectra_assay_overview.pdf}
    \caption[Instrument-1 spectra and assay distributions]{Instrument-1 spectra and supplied assay distributions.  The
    spectral panel gives the median and 10th--90th percentile band for the
    development and designated-benchmark samples; the dashed line marks the
    prespecified 1798~nm boundary used only in a sensitivity analysis.  The
    response panel separates the calibration, validation, and benchmark
    splits.}
    \label{fig:spectroscopy-overview}
\end{figure}

This is an empirically ill-posed functional regression problem.  There are
650 grid coordinates but only 195 development observations, and sample
centring limits the fitted covariance rank to at most 194.  At least 456 grid
directions are consequently unidentified from the development sample.  The
equal-grid convention from Section~\ref{sec:methods-empirical-operator} is
retained:
\begin{equation}
    \langle f,g\rangle_T=T^{-1}f^\top g,
    \qquad
    \widehat K=\frac{X^\top X}{nT},
    \qquad
    \widehat r=\frac{X^\top y}{n}.
    \label{eq:spectroscopy-grid-geometry}
\end{equation}
The resulting prediction for a new spectrum $x$ is
\begin{equation}
    \widehat Y(x)=\bar Y_{\mathrm{train}}
      +\frac{1}{T}(x-\bar X_{\mathrm{train}})^\top\widehat\beta .
    \label{eq:spectroscopy-prediction}
\end{equation}
This preserves the scale of the existing implementation and makes the
training-sample intercept explicit.

\section{Frozen evaluation protocol}
\label{sec:spectroscopy-protocol}

The primary spectra are otherwise unprocessed.  Every wavelength and the
response are centred within each training fold, and those training means are
then applied unchanged to the corresponding validation observations.  There
is no primary variance scaling, smoothing, differentiation, standard-normal-
variate transformation, multiplicative-scatter correction, baseline
correction, or wavelength selection.  This restrained primary analysis
keeps preprocessing from becoming an unreported regulariser.  Prespecified
trimmed-spectrum, smoothed-spectrum, and instrument-2 analyses are reported as
sensitivities in Appendix~\ref{app:spectroscopy-supplement}.

The complete procedures retain their native regularisation rules.  Raw and
Arnoldi FPLS independently minimise prediction mean squared error over at most
70 requested Krylov dimensions, using the same five inner folds.  Raw FPLS
uses unscaled covariance powers and response-space normal equations; Arnoldi
uses twice reorthogonalised Gram--Schmidt and response-space QR.  FPCR uses
response-space GCV over its spectral cut-offs.  CG--FPLS uses the released-code
conjugate-residual recurrence and discrepancy rule with $\tau=1.01$ and
$\delta=0.10$, including its internal moment-GCV variance pilot.  Component
counts therefore do not represent common effective degrees of freedom.  The
comparison concerns four complete regularisation procedures, consistent with
the three-axis framework in Section~\ref{sec:methods-comparison-framework}.

Development performance is assessed by five repeats of five-fold outer
cross-validation.  Folds are constructed separately within the supplied
calibration and validation rows and balanced over assay, so every outer fold
contains 31 calibration and 8 validation tablets.  Raw and Arnoldi receive a
new common five-fold inner partition inside each outer training sample.  The
five repeat-level out-of-fold summaries describe development behaviour; the
25 overlapping fits are not treated as independent replications.

After development checks and code review, each method is fitted once to all
195 development tablets.  A single frozen, common inner partition tunes Raw
and Arnoldi, and the resulting models predict the 460 benchmark tablets once.
The primary endpoint is RMSE.  MAE, signed bias (prediction minus observation),
and $R^2=1-\mathrm{SSE}/\mathrm{SST}$ are secondary.  Conditional uncertainty
is summarised by 10,000 paired iid bootstrap resamples of the benchmark
tablets.  The same resampled tablet indices are used for every method, but the
models are not refitted.  Circular moving-block bootstraps of lengths 5, 10,
and 20 assess possible row-order dependence.  Exact folds and bootstrap
indices are retained in the frozen archive.

\section{Benchmark prediction}
\label{sec:spectroscopy-benchmark}

All four final fits and all benchmark predictions are finite.  Table
\ref{tab:spectroscopy-primary-results} reports the prespecified metrics.
Every method explains substantial variation in the supplied benchmark assay,
with $R^2$ between 0.887 and 0.914.  Raw FPLS has the smallest point RMSE,
4.607, followed by Arnoldi FPLS at 4.643.  FPCR and CG--FPLS have RMSEs of
5.185 and 5.293, respectively.

\begin{table}[htbp]
    \centering
    \small
    \caption[Primary designated-benchmark results]{Primary designated-benchmark results for full-spectrum
    instrument~1 ($n=460$).  Brackets give 95\% percentile intervals from
    10,000 paired iid bootstrap resamples of the benchmark tablets.  The
    intervals shown beside RMSE are marginal method intervals, whereas method
    comparisons use paired contrasts in Table
    \ref{tab:spectroscopy-paired-contrasts}.}
    \label{tab:spectroscopy-primary-results}
    \begin{tabular}{@{}lrrrrr@{}}
        \toprule
        Method & $m$ & RMSE [95\% interval] & MAE & Bias & $R^2$ \\
        \midrule
        CG--FPLS     & 3  & 5.293 [4.367, 6.477] & 3.673 &  0.177 & 0.887 \\
        Raw FPLS     & 55 & 4.607 [3.645, 5.785] & 3.036 & -0.303 & 0.914 \\
        Arnoldi FPLS & 5  & 4.643 [3.693, 5.805] & 3.086 & -0.301 & 0.913 \\
        FPCR         & 53 & 5.185 [4.299, 6.255] & 3.547 & -0.628 & 0.892 \\
        \bottomrule
    \end{tabular}
\end{table}

\begin{figure}[htbp]
    \centering
    \includegraphics[width=0.86\textwidth]{tablet_benchmark_observed_vs_predicted.pdf}
    \caption[Observed and predicted benchmark assay]{Observed and predicted assay on the designated benchmark.  The
    diagonal is the one-to-one line.  The four methods capture the main linear
    signal and share several difficult tablets near the extremes of the assay
    range.}
    \label{fig:spectroscopy-observed-predicted}
\end{figure}

Figure~\ref{fig:spectroscopy-observed-predicted} shows that the four methods
largely make the same errors.  The raw and Arnoldi predictions are especially
close: their correlation is 0.999934, the root mean squared difference between
paired predictions is 0.174 assay units, and the maximum difference is 0.445.
This pairing explains why the two wide marginal RMSE intervals in Figure
\ref{fig:spectroscopy-rmse-intervals} can overlap almost completely while the
paired contrast is estimated much more precisely.

\begin{figure}[htbp]
    \centering
    \includegraphics[width=0.78\textwidth]{tablet_benchmark_rmse_intervals.pdf}
    \caption[Benchmark RMSE estimates and bootstrap intervals]{Benchmark RMSE point estimates and marginal 95\% percentile
    intervals from the paired iid bootstrap.  Overlap of these marginal
    intervals is not a test of a paired method contrast; the declared paired
    differences are reported in Table
    \ref{tab:spectroscopy-paired-contrasts}.}
    \label{fig:spectroscopy-rmse-intervals}
\end{figure}

The focal paired difference in Table
\ref{tab:spectroscopy-paired-contrasts} is
$\mathrm{RMSE}_{\mathrm{raw}}-\mathrm{RMSE}_{\mathrm{Arnoldi}}=-0.036$ assay
units.  Its iid percentile interval excludes zero, as do all three
moving-block sensitivity intervals.  The numerical difference is nevertheless
only 0.78\% of the Arnoldi RMSE.  It is small relative to the benchmark
prediction errors; its practical importance cannot be determined without an
assay scale and an application-specific tolerance.  The interval describes a
conditional difference between two already fitted models on these tablets and
does not establish general superiority of the raw representation.  Arnoldi's
advantages over CG--FPLS and FPCR are larger, approximately 0.65 and 0.54
RMSE units, and also persist under the principal block-length-10 bootstrap.

\begin{table}[htbp]
    \centering
    \small
    \caption[Declared paired benchmark RMSE contrasts]{Declared benchmark RMSE contrasts.  Positive values favour
    Arnoldi FPLS.  The iid and circular block-length-10 columns give 95\%
    percentile intervals from 10,000 paired resamples.  Models are held fixed
    throughout the bootstrap.}
    \label{tab:spectroscopy-paired-contrasts}
    \begin{tabular}{@{}lrrr@{}}
        \toprule
        Contrast & Difference & Iid interval & Block-10 interval \\
        \midrule
        Raw $-$ Arnoldi & -0.036 & [-0.057, -0.017] & [-0.062, -0.014] \\
        CG $-$ Arnoldi  &  0.650 & [ 0.397,  0.921] & [ 0.257,  1.017] \\
        FPCR $-$ Arnoldi&  0.542 & [ 0.286,  0.844] & [ 0.201,  0.922] \\
        \bottomrule
    \end{tabular}
\end{table}

\section{Selection and numerical stability}
\label{sec:spectroscopy-stability}

Prediction alone obscures a much sharper difference between the two leading
procedures.  Across the 25 development outer fits, Arnoldi selects between
four and six components, with median five and no numerical flag.  Raw FPLS
selects between four and 70 requested components, with median 44, reaches the
cap twice, and is flagged in 24 of 25 fits.  Its mean repeated-CV RMSE is also
more variable across the five repeats: $5.082\pm0.517$, compared with
$4.917\pm0.131$ for Arnoldi.  One adverse raw repeat accounts for much of that
difference, and a contiguous-fold sensitivity reverses the leading pair.
These dependent CV summaries therefore diagnose procedure stability rather
than establish another inferential ranking.

\begin{figure}[htbp]
    \centering
    \includegraphics[width=0.82\textwidth]{tablet_development_selected_components.pdf}
    \caption[Development selection and stopping indices]{Selection and stopping indices across the 25 development outer
    fits.  Points show individual fits and boxes summarise their distributions.
    Raw and Arnoldi use inner cross-validation, FPCR uses GCV, and CG--FPLS
    uses discrepancy stopping; counts are therefore not common effective
    degrees of freedom.}
    \label{fig:spectroscopy-selected-components}
\end{figure}

The final raw fit selects nominal dimension 55.  At that point both its
normalised power-basis condition number and its response-design condition
number are recorded at the reporting ceiling, $4.50\times10^{15}$, and the fit
is flagged.  The fixed $1/\sqrt{\epsilon}$ threshold is already crossed by
$m=6$.  Finite predictions do not make those raw coordinates numerically
trustworthy.  Arnoldi selects five components, has basis condition number one,
response-design condition number 11.95, and orthogonality defect
$7.4\times10^{-16}$.  The central empirical finding is consequently not that
one representation wins on predictive loss: Arnoldi realises essentially the
same predictive relation as Raw through stable five-dimensional coordinates.
The complete paths are shown in Appendix~\ref{app:spectroscopy-supplement}.

CG--FPLS stops at three components in every development fit and in the final
fit.  This stable but strongly regularising behaviour accompanies its larger
prediction error.  The finding concerns the implemented discrepancy rule,
including its variance pilot, rather than a general inferiority of conjugate-
gradient regularisation.  FPCR selects between seven and 57 components in
development and 53 in the final fit, illustrating again that nominal component
counts cannot be compared mechanically across different search spaces and
tuning criteria.

\section{Robustness, interpretation, and scope}
\label{sec:spectroscopy-scope}

The prespecified sensitivities confirm that the detailed ranking is conditional
on the measurement and preprocessing choices.  Truncating instrument~1 at
1798~nm reduces the RMSEs of CG--FPLS, Arnoldi, and FPCR and compresses all
four results to 4.63--4.77; Arnoldi then has the smallest point RMSE.  Fixed
Savitzky--Golay smoothing favours Raw at 4.66 but does not improve every
procedure.  Full-spectrum instrument~2 produces larger RMSEs for all four
methods, with Raw and Arnoldi again the leading pair.  Instrument~2 is a paired
robustness analysis of the same tablets, not an external replication or a
calibration-transfer experiment.  These patterns show preprocessing and
instrument dependence, but do not establish that the visible high-wavelength
spikes caused particular prediction errors.

The residual displays in Appendix~\ref{app:spectroscopy-supplement} contain a
few large, common errors and regression toward the centre of the assay range.
All such tablets remain in the analysis.  Without batch, acquisition-order, or
laboratory metadata, row-order patterns cannot be assigned a physical cause.
Likewise, the fitted coefficient vectors identify neither a unique functional
slope nor chemically causal wavelengths in this $650$-coordinate,
rank-$194$ problem.

Three limitations govern the result.  First, the designated benchmark had
been inspected during dataset selection, so the reported errors are not a
prospectively blind validation.  Second, paired-bootstrap intervals condition
on the fitted models and the supplied benchmark tablets; they omit development-
sample, preprocessing, and tuning uncertainty.  Third, the study contains one
response from one tablet archive and no independent batch, laboratory, or
acquisition-period validation.  Subject to those limits, it transfers two
mechanisms from the controlled experiments to an observed problem: aggressive
early stopping can sacrifice prediction, and algebraically related Krylov
fits can have almost identical predictions while their finite-precision
representations have radically different reliability.
